\section{Experiment Design}
\label{sec:design}

\subsection{Datasets}
We evaluate the implementations on five real-world datasets and one synthetic
dataset, summarized in Table~\ref{tab:datasets}. All attributes are numeric.

SUSY and HIGGS~\cite{baldi2014higgs} contain simulated particle collision events
labelled as signal or background. Following the UCI repository, the last 500,000
samples of each serve as the test set. ECBDL'14~\cite{triguero2015ecbdl}
originates from the protein structure prediction task of the ECBDL'14 Big Data
Competition; its nominal attributes are encoded as integers.
MNIST8M~\cite{loosli2007mnist8m} consists of 8.1 million handwritten digits
generated by deforming the MNIST training images, and
SOREL-20M~\cite{harang2020sorel} is a malware detection benchmark built from
Windows executables. These three datasets are used with their official test
sets.

The training sets of ECBDL'14 and SOREL-20M do not fit
into the 80\,GB memory of our GPUs. To include them in the GPU experiments, we
train on their first 15 million and 4 million samples, respectively, and
evaluate the trees on the full test sets.

Chess is a synthetic $3 \times 3$ checkerboard: two attributes drawn uniformly
from $[0, 3]$ and two classes alternating between the unit squares, without
noise. Its exact tree has only nine leaves, so the training time reflects the
cost per sample rather than tree growth. We use sizes from $10^6$ to
$2 \cdot 10^9$ samples, for timing only.

\begin{table}
\caption{Datasets. Size: the training set as float32.}
\label{tab:datasets}
\centering
\small
\begin{tabular}{@{}lrrrrr@{}}
\toprule
Dataset & Attributes & Classes & Training samples & Test samples & Size \\
\midrule
SUSY~\cite{baldi2014higgs}          & 18   & 2  & 4,500,000  & 500,000   & 0.3\,GB \\
HIGGS~\cite{baldi2014higgs}         & 28   & 2  & 10,500,000 & 500,000   & 1.2\,GB \\
ECBDL'14~\cite{triguero2015ecbdl}   & 631  & 2  & 31,992,921 & 2,897,917 & 80.8\,GB \\
\quad GPU subsample                 & 631  & 2  & 15,000,000 & 2,897,917 & 37.9\,GB \\
MNIST8M~\cite{loosli2007mnist8m}    & 784  & 10 & 8,100,000  & 10,000    & 25.4\,GB \\
SOREL-20M~\cite{harang2020sorel}    & 2381 & 2  & 12,699,013 & 4,195,042 & 120.9\,GB \\
\quad GPU subsample                 & 2381 & 2  & 4,000,000  & 4,195,042 & 38.1\,GB \\
Chess 3$\times$3                    & 2    & 2  & $10^6$--$2 \cdot 10^9$ & \multicolumn{1}{c}{--} & 8\,MB--16\,GB \\
\bottomrule
\end{tabular}
\end{table}

\subsection{Hardware}
\plan{Four nodes of the mordor3 cluster (Table~\ref{tab:hardware}), each used
whole (exclusive jobs). GPU nodes run the CUDA backend only, CPU nodes the CPU
backends and the other CPU implementations. GPU runs are bound to the GPU's
NUMA node (16 CPUs on hopper, 24 on ampere). Add measured memory bandwidth
(STREAM) per node. Software: Ubuntu 24.04, Linux 7.0, g++ 13.3 (\texttt{-O3}),
CUDA 12.9, NVIDIA driver 615.71.}

\begin{table}
\caption{Hardware. Every node has two sockets with SMT disabled; CPU, cores and L3
are per socket. Clock is base/max turbo in GHz. One GPU is used per run.}
\label{tab:hardware}
\centering
\small
\setlength{\tabcolsep}{5pt}
\begin{tabular}{@{}llrrrll@{}}
\toprule
Node & CPU & Cores & Clock & L3 & RAM & GPU \\
\midrule
hopper & Xeon Gold 6444Y & 16 & 3.6/4.0 & 45\,MB  & 1\,TB DDR5-4800   & H100 80\,GB \\
ampere & EPYC 7402       & 24 & 2.8/3.35& 128\,MB & 1\,TB DDR4-3200   & A100 80\,GB \\
epyc  & EPYC 9255       & 24 & 3.25/4.3& 128\,MB & 512\,GB DDR5-6400 & -- \\
dell  & Xeon Gold 6542Y & 24 & 2.9/4.1 & 60\,MB  & 1\,TB DDR5-5200   & -- \\
\bottomrule
\end{tabular}
\end{table}

\subsection{Compared Implementations}
\plan{./tree backends (\texttt{--simple}, \texttt{--serial},
\texttt{--parallel}, \texttt{--cuda}, \texttt{--gpu-only}); scikit-learn and
Weka J48 on the CPU; cuML and XGBoost on the GPU, with their settings. Why each
is the closest equivalent (no public GPU C4.5; J48 depth limit adapter;
XGBoost squared error = Gini on bins). rpart and YaDT: add or say why not.}

\subsection{Protocol and Measurement}
\plan{One tree, depth $\le$ 10, no pruning, all samples and attributes; what is
timed (training only, from data in memory; GPU setup and upload included);
peak RAM and VRAM; runs per case, means and spread; exclusive nodes; fairness
(CPU threads of the hybrid, warm library processes, input layout). Source:
docs/GPU-LIBRARIES.md, bench/README.md.}
